% =============================================================================
% Certus -- Section: Method (DRAFT, Week 8)
% Code references are to src/certus/; every number in Setup/Results is
% regenerable from scripts/ and results/.
% =============================================================================

\section{Method}
\label{sec:method}

Certus is a calibration layer, not a detector. It takes any model that
assigns an image a scalar score, uses a labelled calibration set to turn
that score into a triage threshold, and attaches a finite-sample,
distribution-free bound on the rate at which real defects are silently
cleared. The one assumption is that calibration and deployment images are
exchangeable; \S\ref{sec:label-noise} and the shift study test what
happens when they are not. The detector is never retrained. Everything below holds for
\emph{any} score; the score only determines how many images must go to a
human.

\subsection{Scores and the triage rule}

A nonconformity score $s : \mathcal{X} \to \mathbb{R}$ maps an image to a
real number. We orient every score the same way: a high $s(x)$ means the
model sees little evidence of a defect. The triage rule has one parameter
$\lambda$:
\begin{equation}
  \textsc{auto-clear}(x) \iff s(x) > \lambda,
\end{equation}
and every image not auto-cleared is routed to a human. The economic output
of the system is the \emph{human-review fraction},
$\mathrm{HRF}(\lambda) = \Pr(s(X) \le \lambda)$.

\subsection{Missed-defect loss and two guarantees}

With $y \in \{0, 1\}$ indicating a defect, the missed-defect loss is
\begin{equation}
  L_i(\lambda) = \mathbf{1}\{y_i = 1 \wedge s(x_i) > \lambda\},
\end{equation}
the indicator that a real defect was auto-cleared. $L_i$ is non-increasing
in $\lambda$: raising the threshold can only clear fewer images. This
monotonicity is the one property the guarantees below rely on. The code
checks it at every calibration, and a unit test confirms that the
sign-flipped rule is rejected rather than silently accepted.

We report two guarantees that differ in what they condition on.
\begin{itemize}
  \item \textbf{Joint:} $\mathbb{E}[\mathbf{1}\{\text{defect} \wedge
  \text{cleared}\}] \le \varepsilon$, the share of \emph{all} images that
  are defects and were cleared. Calibrated on all rows.
  \item \textbf{Conditional:} $\Pr(\text{cleared} \mid \text{defect}) \le
  \varepsilon$, the share of \emph{real defects} that are missed.
  Calibrated on defect-positive rows only, which are exchangeable with
  future defects. This is the statement an inspector actually wants. It is
  more expensive, because the defects alone must certify the bound.
\end{itemize}

\subsection{Conformal risk control}

Given $n$ exchangeable calibration examples, conformal risk control
\citep{angelopoulos2024crc} sets
\begin{equation}
  \hat\lambda = \inf\Big\{\lambda : \tfrac{n}{n+1}\hat R_n(\lambda) +
  \tfrac{B}{n+1} \le \varepsilon\Big\},
  \qquad \hat R_n(\lambda) = \tfrac{1}{n}\textstyle\sum_i L_i(\lambda),
\end{equation}
with $B = 1$. This gives $\mathbb{E}[L_{n+1}(\hat\lambda)] \le \varepsilon$
over the joint draw of calibration and test data, and the bound is tight to
$O(1/n)$. Because $\hat R_n$ only changes at calibration scores, the
infimum is searched over those scores plus $+\infty$, the
``clear nothing'' fallback. That fallback is returned whenever no finite
threshold passes, which is guaranteed once $\varepsilon < B/(n+1)$. In that
case the honest answer is that every
image goes to a human.

\subsection{Learn-then-Test}

CRC bounds risk \emph{on average} over calibration draws, so a single
deployment can land above $\varepsilon$: in our runs 31--57\% of seeds
did, though by small margins.
Learn-then-Test \citep{angelopoulos2021ltt} instead gives
$\Pr(R(\hat\lambda) \le \varepsilon) \ge 1 - \delta$ for \emph{this}
deployment. For each $\lambda$ it tests $H_\lambda : R(\lambda) >
\varepsilon$ with the Hoeffding--Bentkus p-value
\begin{equation}
  p_\lambda = \min\Big(e^{-n\,h(\hat R_n \wedge \varepsilon,\,
  \varepsilon)},\; e\cdot\Pr\big(\mathrm{Bin}(n, \varepsilon) \le \lceil
  n\hat R_n\rceil\big)\Big),
\end{equation}
where $h$ is the Bernoulli KL divergence. It then applies fixed-sequence
testing from $\lambda = +\infty$ downwards, stopping at the first
non-rejection. Monotone risk makes this order natural, so $\delta$ is
not split across hypotheses.

\subsection{Split conformal prediction}

For the validity study (E1) and as baselines we also implement split
conformal prediction with the LAC, APS and RAPS scores. The threshold
$\hat q$ is the $\lceil(n+1)(1-\alpha)\rceil$-th smallest calibration
score, computed as an exact order statistic. The common
$\mathrm{quantile}(1-\alpha)$ shortcut interpolates below that order
statistic and under-covers. Our property test catches it at
$n = 20$, $\alpha = 0.1$, where expected coverage drops to about $0.86$.

\subsection{Nonconformity scores}

No fine-tuned detector exists for our data: the source dataset's public
project lists zero trained models. Every score is therefore built from
off-the-shelf models, which the method's model-agnosticism permits and
which makes the comparison reproducible.
\begin{description}
  \item[S1, CLIP zero-shot.] $1 - p(\text{defect})$, where $p$ is the
  CLIP ViT-B/32 softmax mass on four defect prompts against two clean-wall
  prompts. A linear head trained on the bootstrapped labels was at chance
  (5-fold AUC $0.51$); zero-shot needs no labels at all.
  \item[S2, generative log-probability.] $-\log p(\text{``Yes''})$ as
  Qwen2-VL-2B's answer to a fixed defect question, teacher-forced.
  Averaging over a full assertion sentence diluted the one informative
  token (AUC $0.58$ vs.\ $0.82$ on an 80-image selection sample). That
  single binary choice is the only one made on calibration-split data:
  31 of the 80 selection images fall in the 3000-image evaluation pool
  (1\%). Score quality is also reported on the disjoint tuning set.
  \item[S3, verbalised confidence.] $1 - c/100$ for the model's stated
  0--100 confidence. Included as the practitioner baseline.
  \item[S4, detector objectness.] $1 - $ the maximum OWLv2 box confidence
  over one text query per defect class.
  \item[S5, MC-dropout bound.] Twenty passes of dropout ($p = 0.1$) on the
  CLIP embedding before S1's prompt comparison. The score is
  $\mathrm{mean}(1-p_k) - \mathrm{std}(1-p_k)$: an image is cleared only if
  it looks clean \emph{and} stays clean under perturbation. Masks are
  seeded per image.
  \item[S6, feature typicality.] Negative Mahalanobis distance (Ledoit--Wolf
  covariance) from the CLIP embedding to 3000 defect-free training images.
  Defects and unfamiliar scenes both score low.
  \item[S7, learned combination.] Logistic regression over standardised
  S1--S6.
\end{description}

\subsection{Three splits, not two}

A score fitted on calibration data is no longer independent of it, and the
guarantee silently fails. Both fitted scores therefore learn only from
the training split: S6's reference set is 3000 clean training images, and
S7's combiner is fitted on a separate 1000-image tuning set. Neither
overlaps any calibration pool. This is enforced in code: every
function whose name begins with \texttt{fit\_} calls a guard that raises if
it is handed a calibration or test row, and S7 refuses to score its own
tuning rows.
